% Generated draft: requires mathematical and source review.
\documentclass[11pt]{article}
\usepackage[margin=25mm]{geometry}
\usepackage{fontspec}
\setmainfont{Latin Modern Roman}
\usepackage{amsmath,amssymb,mathtools,unicode-math}
\setmathfont{Latin Modern Math}
\usepackage{enumitem,needspace,xurl,hyperref}
\setlength{\parindent}{0pt}
\setlength{\parskip}{5pt}
\setlength{\emergencystretch}{3em}
\newcommand{\N}{\mathbb{N}}
\newcommand{\Z}{\mathbb{Z}}
\newcommand{\Q}{\mathbb{Q}}
\newcommand{\R}{\mathbb{R}}
\newcommand{\C}{\mathbb{C}}
\newcommand{\F}{\mathbb{F}}
\newcommand{\A}{\mathbb{A}}
\newcommand{\PP}{\mathbb{P}}
\newcommand{\E}{\mathbb{E}}
\newcommand{\eps}{\varepsilon}
\newcommand{\dd}{\,\mathrm{d}}
\begin{document}
\section*{UnsolvedMath ID: 801080}
\textbf{Source problem code:} catalogue entry
\subsection{Definitions and mathematical statement}
A standard convention for a generalized H\'enon map is a finite composition of maps
\[
H_j(x,y)=(y,p_j(y)-a_jx),\qquad p_j\in K[t],\quad \deg p_j\geq 2,\quad a_j\in K\setminus\{0\},
\]
acting on \(K^2\), where \(K=\R\) or \(K=\C\). This convention explains the terminology; the supplied excerpt does not specify the ground field, the precise class of maps, or any restrictions on their parameters, so these remain unresolved hypotheses.

If \(z\) is a periodic point of a differentiable map \(f\) with least period \(n\geq1\), its multipliers are the eigenvalues of \(D(f^n)_z\), taken over \(\C\). An unstable multiplier is an eigenvalue \(\lambda\) satisfying \(|\lambda|>1\). At a saddle periodic point of a two-dimensional map, the two eigenvalues satisfy
\[
|\lambda_s|<1<|\lambda_u|,
\]
and \(\lambda_u\) is its unique unstable multiplier. The excerpt does not specify whether its list includes only saddle periodic orbits or a different collection of periodic points.

Once the map class, the periodic objects included, and the meaning of a list have been fixed, let \(\mathcal U(f)\) denote the resulting unstable-multiplier data and let \(\mathcal U_{\mathrm{abs}}(f)\) denote the corresponding data obtained by replacing each multiplier \(\lambda\) by \(|\lambda|\), retaining whatever indexing and multiplicity convention was chosen.

\textbf{Question 32 (Berger and Bedford, 2023).} To what extent does \(\mathcal U(f)\), respectively \(\mathcal U_{\mathrm{abs}}(f)\), determine a generalized H\'enon map \(f\)? Equivalently, characterize the possible maps \(g\) satisfying, respectively,
\[
\mathcal U(g)=\mathcal U(f)
\qquad\text{or}\qquad
\mathcal U_{\mathrm{abs}}(g)=\mathcal U_{\mathrm{abs}}(f).
\]
The intended notion of determination is not specified: the excerpt does not state an equivalence relation on maps or a proposed bound on the remaining freedom. It also does not specify whether periods, orbit labels, or multiplicities belong to the spectral data. Consequently, the supplied question does not determine a unique rigidity assertion without further source clarification.

The source motivates this question by saying that a local unstable slice contains essentially complete information about unstable multipliers and by relating it to the preceding question and classical spectral rigidity. The preceding question and the meaning of this qualification are not supplied; this motivation is not asserted here as a precise theorem.
\subsection{Short English statement}
Berger and Bedford ask how much a generalized H\'enon map is determined by its unstable multipliers, or just their absolute values. A precise rigidity statement still requires the map class, the periodic data and their indexing, and the intended equivalence of maps to be specified. The supplied registry labels the problem open; this draft does not independently verify that status.
\subsection{Sources}
\begin{itemize}
\item Local UnsolvedMath registry, permanent ID \(801080\), source code \texttt{catalogue entry}, title \emph{Henon Maps: Open Problems --- Question 32}. The supplied record attributes the question to Berger and Bedford, with proposed year 2023, and labels it \texttt{open}.
\item Berger and Bedford, \emph{Henon Maps: A List of Open Problems} (2023), Question 32, \url{https://arxiv.org/abs/2312.03907}. This reference and attribution are supplied by the registry; the referenced document has not been independently consulted for this draft.
\item Supplied provenance: access date 2026-10-16; extraction method \texttt{arxiv\_latex}; upstream revision \texttt{d371ee917ac9647fbb06cd6fc529deec6f3dda82}. The access date records supplied provenance, not an independent source verification or literature assessment.
\end{itemize}
\end{document}
